\section{Conclusion}
\label{sec:conclusion}

We introduced Minimum Intervention Reweighting, a dual-constraint framework for contamination-aware pretraining that selects per-batch exponents $b^*$ without validation data.
The method finds the smallest $b$ satisfying $\alpha_{\mathrm{eff}}(b) \geq \tau$ subject to $\mathrm{ESS}(b) \geq n_{\min}$, enforcing a minimum-intervention principle: apply only as much downweighting as needed to meet a quality floor, no more.
This stands in direct contrast to ESS-only (maximum-intervention) strategies, which we show degrade performance below even uniform baselines by over-downweighting contaminated samples.

Across three seeds at 0.6B scale with 40\% synthetic contamination, adaptive $b^*$ matches grid-search-optimal PPL (17.651 vs.\ 17.646) within seed variance, eliminating the $|\mathcal{B}|$-run grid search entirely.
The quality threshold $\tau$ is robust across $[0.5, 0.8]$---all values in this range outperform uniform weighting.
A lightweight TF-IDF + logistic regression detector (AUC\,0.94) suffices when post-hoc calibrated (ECE\,0.048); score calibration matters more than raw discriminative power.
The batch-level constraint regime---49.9\% clean (no reweighting needed), 20.2\% $\alpha_{\mathrm{eff}}$-binding, 30.0\% ESS-binding---confirms that both constraints are active and complementary in practice.

\section*{Limitations}
\label{sec:limitations}

\paragraph{Scale.}
All experiments use a single model size (Qwen3-0.6B) with batch size 4.
Batch-level $q$-score statistics concentrate differently at larger batch sizes typical of 3B+ training; the optimal $\tau$ range and constraint dynamics may shift at scale.

\paragraph{Contamination distribution.}
We inject synthetic data via uniform random mixing at a controlled ratio ($\rho{=}0.4$).
Real web-scale contamination is domain-skewed, which may alter within-batch $\alpha_{\mathrm{eff}}$ distributions and require domain-aware extensions.

\paragraph{Detector coverage.}
We evaluate one detector family (TF-IDF + logistic regression).
Although our results suggest calibration dominates discrimination for reweighting quality, a systematic ablation across detector architectures and quality levels is needed.

\paragraph{Threshold.}
$\tau$ is an empirical hyperparameter.
The convergence theory of \citet{wang_2026_survive_contamination} assumes a known contamination proportion $\alpha$, which we do not require; bridging this gap to derive $\tau$ from first principles remains open.

\paragraph{Evaluation.}
We report held-out perplexity only.
Downstream task evaluation (MMLU, HellaSwag) would strengthen the practical case, particularly at scales where PPL differences translate more reliably to task performance.

\paragraph{Effect size.}
The PPL improvement over uniform weighting is 0.043 at 0.6B/40\% contamination.
While consistent across seeds, this effect is modest; we expect amplification at higher contamination ratios and larger model scales, but verification is ongoing.
